%% kapa: exam=1 round=34 year=2023 period=1 form=A subject=real_estate subject_ko=부동산학원론 seq=39 number=119
%% source: Q-Net 감정평가사 기출문제 (cst00309) artlSeq=5211789 file=raw/1008/5211789/2023년도 제34회 감정평가사 1차시험 1교시.pdf
%% conversion: kordoc 4.24.51 -> tools/gen1.py (auto); review=proofread-sonnet; flags=none
\begin{question}{119}
다음 자료를 활용하여 공시지가기준법으로 평가한 대상토지의 시산가액(m$^2$당 단가)은?
\begin{framedtext}
○ 대상토지 현황: A시 B구 C동 101번지, 일반상업지역, 상업나지 \par
○ 기준시점: 2023.04.08. \par
○ 비교표준지: A시 B구 C동 103번지, 일반상업지역, 상업나지 \par
\hspace*{4.5em}2023.01.01.기준 표준지공시지가 10,000,000원/m$^2$ \par
○ 지가변동률: 1) 2023.01.01. \~{} 2023.03.31. : -5.00 \% \par
\hspace*{6.2em}2) 2023.04.01. \~{} 2023.04.08. : -2.00 \% \par
○ 지역요인: 비교표준지는 대상토지의 인근지역에 위치함 \par
○ 개별요인: 대상토지는 비교표준지대비 획지조건에서 4\% 열세하고, 환경조건에서 5\% 우세하며, 다른 조건은 동일함 \par
○ 그 밖의 요인 보정: 대상토지 인근지역의 가치형성요인이 유사한 정상적인 거래사례 및 평가사례 등을 고려하여 그 밖의 요인으로 20\% 증액 보정함 \par
○ 상승식으로 계산할 것 \par
○ 산정된 시산가액의 천원미만은 버릴 것 \par
\end{framedtext}
\begin{choices}
  \item 11,144,000원
  \item 11,168,000원
  \item 11,190,000원
  \item 11,261,000원
  \item 11,970,000원
\end{choices}
\end{question}
\answer{4}
